\documentclass{article}

\usepackage{booktabs}
\usepackage{tabularx}

\title{Development Plan\\\progname}

\author{\authname}

\date{}

\input{../Comments.text}
\input{../Common.text}

\begin{document}

\maketitle

\begin{table}[hp]
\caption{Revision History} \label{TblRevisionHistory}
\begin{tabularx}{\textwidth}{llX}
\toprule
\textbf{Date} & \textbf{Developer(s)} & \textbf{Change}\\
\midrule
Date1 & Name(s) & Description of changes\\
Date2 & Name(s) & Description of changes\\
... & ... & ...\\
\bottomrule
\end{tabularx}
\end{table}

\newpage{}

The purpose of this document is to outline the development plan for Magic 8, a solitaire game focused on finding missing numbers on an electronic 8x8x8 magic cube. This will contain information on how the team is set up and will meet, the workflow they will be using, alongside their PoC Demonstration Plan and their AI usage.


\section{Confidential Information?}

There is no confidential information to protect as this project will be based on existing, open access literature. 

\section{IP to Protect}

There is no IP to protect, and the team has agreed to let Dr. Paul Rapoport post it for free on a website of his choosing.

\section{Copyright License}

This project is licensed under the MIT License. The full license can be found at \href{https://github.com/Jacob-MB/magic-8/blob/main/LICENSE}{magic-8/LICENSE}.

\section{Team Meeting Plan}

Circle Back to This

\wss{This section is for meetings between the team members, and meetings between 
the team and their supervisor/advisor (as appropriate).}

\wss{How often will the team meet? where? when?}

\wss{If the meeting is a physical location (not virtual), out of an abundance of
caution for safety reasons you shouldn't put the location online}

\wss{To help schedule meetings, you can use the lecture/tutorial time when we do 
not have anything else scheduled.}

\wss{How often will you meet with your supervisor/advisor?  when?  where?}

\wss{Will meetings be virtual?  At least some meetings should likely be
in-person.}

\wss{How will the meetings be structured?  There should be a chair for all meetings.  
There should be an agenda for all meetings.}

\section{Team Communication Plan}

We will be using the following communication tools and methods to ensure that this is a successful project.

\begin {itemize}
\item \textbf{Github}: There is a Github repository that will contain the code and documentation for this project. There will be two releases: one that will allow you to run the project as an application, and another which will be used for Dr. Rapoport's prescribed website. In lieu of assigning tasks in passing and promptly forgetting, there will be a Github Projects board that will follow the Kanban method taught in 3RA3, with tasks assigned simultaneously during meetings. Meeting minutes will also be uploaded to the repository.
\item \textbf{MS Teams}: There is a Teams chat with all 5 group members. This will be the standard place for communication outside of in-person meetings and can be used as either a place for virtual meetings or as an asynchronous chat. Ideation outside of meetings will occur on an appropiate shared Words document, which will be committed onto Github when appropiate. There is also a seperate shared group chat with Dr. Rapoport, which will be used to schedule meetings and ask quick questions. This can also be used as a place to meet virtually with him, but in-person meetings are preferred. 
\item \textbf{Emails/Google Meet}: When corresponding with outside channels, such as other professors and/or stakeholders, emails and Google Meet are the preferred methods. When emailing, the message should be approved in the Teams chat before being sent, and it should be carbon-copied to all other members and Dr. Rapoport. Google Meet was chosen as stakeholders are lekely tech illiterate and Meet is an accessible browser-based meeting platform. 
\end{itemize}

\section{Team Member Roles}

While the team has different strengths, there will be a lot of learning on the fly done for this project, so the team didn't feel comfortable assigning roles yet. This will be done after further discussion and after everyone has taken a stab at deeply exploring the subject matters at hand (ex. 3D visualization). The load and therefore role might end up falling on multiple shoulders or changed, if we see that it is improbable to be completed by one person. The document will then be revised appropiatedly.

Jacob Broderson will serve as the \textbf{Team Liason}, as he has already kept up communication with the supervisor via emails before moving to the Teams chat, and will also be responsible for communicating directly with the course's professors and teaching assistants should the need arise. He will also be primarily responsible for seeking out stakeholders and scheduling meetings with them.

Here are the following unassigned roles, based off of our approved project description:
\begin {itemize}
\item \textbf{Systems Analyst}: As this project deals with a stringent set of rules that makes the Magic Cube work in octal, this person will be responsible for first compiling and creating all of the rules in our chosen programming language. He is reponsible for ensuring that none of the other members' work break these rules, using test cases and static/dynamic analysis tools. As the design space for octal magic cubes is unfinished and there are only 384 confirmed cubes, it is possible for this member to not only discover new ones, but also make/edit the base rules used to make cubes.
\item \textbf{Game/Creative Director}: Basing their work on the rules set by the Systems Analyst and the model/UI produced by the Front End Developer, the Game/Creative Director will focus on rules for the levels of the game, making sure that they are engaging while still educational. Dr. Rapoport already has an idea for the game in mind, and the Game/Creative Director will work in lockstep with him to find a middleground that is perfect with the stakeholders. This director will work in implementing and testing stakeholder feedback, and will also be responsible for thinking up and integrating new game modes.
\item \textbf{Front End Developer (and 3D Artist)}: They will be responsible for creating the user interface based on what the Game/Creative Director requires. They will also learn about 3D modelling/visualization on the side and create the needed visualization that can be efficiently displayed on a website.
\item \textbf{Back End Developer}: This role requires the group member to store the required data (such as the 384 possible cube formations) in an accessible database, and create the pipeline and interfaces that makes it easier for the other developers to access the information. They can also be responsible for account management (especially important for keeping track of user scores and if we implement an online mode) and other managerial interfaces that simplify the jobs for other developers. If we end up choosing to do a family of games, they will need to come up with a creative way to support multiple databases of different sizes.
\item \textbf{Webmaster}: Based on what Dr. Rapoport wants to do display the information (either as an executable linked on a website, or as a website itself), the Webmaster is responsible for writing and deploying the code that displays the work of the Front and Backend Developer. They will have to ensure that the website is accessible to stakeholders in high school, which means using a public domain instead of using localhost or a private server. 
\end {itemize}



\section{Workflow Plan}
Git will be used to keep track of progress, new features that have been brainstormed, and as a stable workflow where we can always point to what everyone is working on and our latest stable project.

After being assigned a task/issue from the Github Projects Kanban Board, the member will create their own branch from a corresponding primary branch (not main). 
The main branch will always be stable, and almost have same status as a new release (it shouldn't be merged into without being peer reviewed and without the express approval of every member).

When the team starts coding and building documentation after being assigned a role, they will start by creating a primary branch named after their role. For example, the Systems Analyst would call their primary branch, systems-analyst.
Whenever there are tasks that requires multiple team members, they create a temporary primary branch named after the members using the branch. The goal of the users on this branch is to merge this branch into main as soon as possible as the mental calculation is that any feature requiring multiple members is worthy of being a release. 

The team member will be allowed to experiment on their primary branch, as they can essentially create sub branches via "git branch role/experiment-description role" (ex. git branch systemanalyst/experiment-adding\_new\_rules systemananalyst).
The member's primary branch should also be stable and when they get assigned an issue, they will move the issue from the "To Do" column to the "In Progress" column. After this has occurred, the person will work on a subbranch they made on their primary branch. For example, the Back End Developer would use the following command: "git branch backenddeveloper/feature-adding\_magic\_squares backenddeveloper".
This will create a copy of the latest commits on backenddeveloper into the subbranch. 

When they are finished working, they will have to git checkout into the primary branch and then create a merge request between the two branches (ex. going to backenddeveloper and making a merge request between backenddeveloper backenddeveloper/feature-adding\_magic\_squares).
They will then move the issue to the "Done (before Review)" column and add the merge request into the task on the Github Projects Kanban Board. They will also remove the previous tag (ex. remove "feature" or "refactor"), add "review", and assign another team member to the issue. These requects will only require a single peer review before being moved to "Done (after Review)" and being accepted.

There can also be a primary branch for documentation, however these aren't as strict as actual code requests and the only requirement to merge into main is that the primary branch (ex. devplan for this document) is stable. 


In terms of issues outside of the coding, they will also be used to track meetings with the group, professors (during lectures) and our supervisor. There will be a title and a checklist for atendance. Outside of this, it will be a place for people to store their questions that can be updated, so everyone is aware of the questions that will be asked and to prevent forgetfullness.

This is more important for documentation as Github Issues is an easy way to assign parts of the written work to team members. There will be a main issue named the deliverable (ex, Development Plan) with a hard deadline. There will then be titled subissues with a soft deadline and an assigned member (ex. Workflow Plan), and they will include a self-contained checklist with short descriptions (ex. How will you be managing issues, including template issues, issue
classification, etc.?). 

The CI pipeline would consist of the following linter. This was chosen because ... There will also be static tests that were created by the Systems Analyst consistently running on the code when it has been commited alongside domain-specific tests written by the specific "branch manager".
Based on what the team member is writing, there might be specific configurations that the member might use (such as the Front End Developer requiring a plugin to display their model). This means that the JSON/YAML file will also have to build before it can it can be merged. There is also a code review in order to ensure that small problems don't become big problems when it comes time to merge from the primary branches into main.

There is also a CD component if we end up publishing the website earlier to get user data, which will require an appropiate rollback strategy. This will ensure that data leaks (such as user data) are mitigated as soon as possible.

\section{Project Decomposition and Scheduling}

The Github project is located at: 
The Kanban board has the following columns: Lectures/General Ideas, Meetings, To Do, In Progress, Done (before Review), Done (after Review).

The Lecture column will contain labels that describe it as either "Attendance" (to ensure that everyone is able to get the required knowledge to complete milestones and its corresponding documents) and/or "Questions" (if we have questions in the descriptions that we want to ask the professor after class). 
There is also space to create issues that are just spitballed ideas.

Meetings will be labeled to show who it is with, and will be populated with pre-meeting questions. During the meeting, meeting minutes will be added to the desciption, alongside roles/tasks assigned. After the meeting is concluded, it will be labelled as "Not Assigned", until the task discussed in meetings are turned into issues (in "To Do") that are labeled to members.
The rest of the board will follow the guidelines set in Workflow Plan, alongside belonging to the appropiate milestone listed below:


\begin{center}
\begin{tabular}{||c c||} 
 \hline
 Milestones & Week \\ [0.5ex] 
 \hline\hline
 Team Formed, Project Selected & Week 03 \\ 
 \hline
 Problem Statement, POC Plan, Development Plan & Week 04 \\
 \hline
 Req.\ Doc.\ and Hazard Analysis Revision 0 & Week 06 \\
 \hline
 V\&V Plan Revision 0 & Week 08 \\
  \hline
 Design Document Rev -1 and CD-Fall Rev 0 & Week 10 \\
  \hline
 Proof of Concept Demo + Interviews & Week 11+12 \\
  \hline
 Design Document Revision 0 & Week 16 \\
  \hline
 Revision 0 Demo + Interviews & Week 18+19 \\
  \hline
 V\&V Report Rev 0 and CD-Winter Rev 0 & Week 22 \\
  \hline
 Final Demonstration (Revision 1) & Week 24 \\
  \hline
 EXPO Demonstration & Week 26 \\
 \hline
 Final Documentation (Revision 1) & Week 26 \\ [1ex] 
 \hline
\end{tabular}
\end{center}

The Final Documentation will contain the Problem Statement, Development Plan, Proof of Concept (POC) Plan, Requirements Document, Hazard Analysis, Design Document, V\&V Plan, V\&V Report, CD-Fall, CD-Winter and the Source Code.


\section{Proof of Concept Demonstration Plan}

Two of the greatest risks currently identified for this project are the 3D visualization of the magic cube and the implementation of the base design of the game and its eight stages. The 3D visualization represents the largest risk that the POC will attempt to address. 
This is because, for the game to function as currently envisioned, 3D visualization is a hard requirement; however, the group has little to no experience working with 3D visualization. The other risk to be addressed is the base logic of the game itself. 
This has been identified as a risk because if the base functionality cannot be implemented, then any further development would be ineffective.

To address these risks, the POC will consist of a minimal version of the game in which we will be able to demonstrate a user moving through a 3D magic cube and completing the eight tasks outlined in the project description. 
These tasks will be simplified and hardcoded versions of the final product, where the numbers provided to the user will always remain the same. This will allow us to test the functionality of the tasks without requiring the random elements of the game to be implemented. 
The POC will be considered successful if the magic cube can be rendered and interacted with in 3D, and if the user can successfully complete all eight tasks without any critical issues with the underlying game logic.

If the demo were to fail, the consequences for the project as a whole would depend on which risk caused the failure. If the implementation of the eight stages of the game were to fail, a significant discussion regarding the future of the project would be necessary. 
This would indicate that the group is unable to deliver the fundamental functionality required by the project and may require a significant deviation from the original capstone plan. 
If the 3D visualization were to fail, the group would instead need to evaluate the possibility of transitioning to a static or 2D version of the visuals rather than continuing with the current 3D plan.

\section{Expected Technology}

\wss{What programming language or languages do you expect to use?  What external
libraries?  What frameworks?  What technologies.  Are there major components of
the implementation that you expect you will implement, despite the existence of
libraries that provide the required functionality.  For projects with machine
learning, will you use pre-trained models, or be training your own model?  }

\wss{The implementation decisions can, and likely will, change over the course
of the project.  The initial documentation should be written in an abstract way;
it should be agnostic of the implementation choices, unless the implementation
choices are project constraints.  However, recording our initial thoughts on
implementation helps understand the challenge level and feasibility of a
project.  It may also help with early identification of areas where project
members will need to augment their training.}

Topics to discuss include the following:

\begin{itemize}
\item Specific programming language
\item Specific libraries
\item Pre-trained models
\item Specific linter tool (if appropriate)
\item Specific unit testing framework
\item Investigation of code coverage measuring tools
\item Specific plans for Continuous Integration (CI), or an explanation that CI
  is not being done
\item Specific performance measuring tools (like Valgrind), if
  appropriate
\item Tools you will likely be using?
\end{itemize}

\wss{git, GitHub and GitHub projects should be part of your technology.}

\section{Use of AI and LLMs}

The team expects to use AI tools, such as GitHub Copilot and Claude, at various stages of the project during both development and documentation. Most uses will fall under brainstorming, creating boilerplate code, and assisting with debugging.

Since the game's core logic is based on a fixed, formula-driven set of rules, any AI-generated implementation of these rules will not be merged until it has been checked against the correct magic sum (2052 decimal / 4004 octal) using the static test suite described in the Workflow Plan. AI-generated 3D visualization code will similarly be manually tested, since correctness in this area cannot be verified by unit tests alone.

For documentation, AI tools may be used to help draft, restructure, or brainstorm sections of project documents. Any AI-assisted content will be reviewed by the team for accuracy before inclusion, particularly where it touches the mathematical rules of the game or project requirements.

Regardless of the source, all code passes through the same merge request and peer review process described in the Workflow Plan. AI-generated contributions are held to the same standard as human-written ones and do not skip any step of the process.

\section{Coding Standard}

\wss{What coding standard will you adopt?}

\newpage{}

\section{Appendix --- Generative AI Use Disclosure}

\subsection{AI Tools Used}

Claude Sonnet 5 (Medium) by Anthropic

\subsection{How and Where Used}

Claude was used to brainstorm ideas for how the team can effectively use AI and LLMs throughout the project. These ideas were used to help write the Use of AI and LLMs section.

\subsection{Verification of AI-Generated Content}

The team reviewed and edited the AI-generated ideas before including them, and verified that they were consistent with the workflow, testing, and peer review processed described in the Workflow Plan. The team remains responsible for the accuracy and completeness of all content in this document, regardless of whether it was generated by AI or written by a team member.

\newpage{}

\section*{Appendix --- Reflection}

\wss{Not required for CAS 741}

\input{../Reflection.text}

\begin{enumerate}
    \item Why is it important to create a development plan prior to starting the
    project?
    \item In your opinion, what are the advantages and disadvantages of using
    CI/CD?
    \item What disagreements did your group have in this deliverable, if any,
    and how did you resolve them?
\end{enumerate}

\subsection*{Allen Tong}

\begin{enumerate}
  \item A development plan is important to create before starting the project because it gets the team to agree on how we will work before any code is written. Deciding on how we will communicate, manage our code, and track tasks is important to ensure that work is not assigned in passing and forgotten, which helps the project run smoothly. It also made us identify our biggest risks, which are the 3D visualization and the base logic of the game, allowing us to plan our Proof of Concept around them to avoid delays later on. Finally, it gives us a shared reference to check against if someone falls behind or the team disagrees on the direction of the project.
  \item The main advantage of using CI/CD is that every merge request is checked automatically before it is merged into the main branch. This is especially important for our project because the game rules are exact, so a test confirming that every line sums to the correct magic constant will catch any mistakes right away. A disadvantage is that it can be time-consuming to set up and maintain the pipeline, especially since most of us are still learning the tools involved.
  \item Our group did not have any disagreements during this deliverable. We split up the tasks beforehand so that everyone knew what to do and what was expected of them. As a result, the team worked well together and everyone understood what the document needed.
\end{enumerate}

\newpage{}

\section*{Appendix --- Team Charter}

\wss{borrows from
\href{https://engineering.up.edu/industry_partnerships/files/team-charter.pdf}
{University of Portland Team Charter}}

\subsection*{External Goals}

\wss{What are your team's external goals for this project? These are not the
goals related to the functionality or quality fo the project.  These are the
goals on what the team wishes to achieve with the project.  Potential goals are
to win a prize at the Capstone EXPO, or to have something to talk about in
interviews, or to get an A+, etc.}

\subsection*{Attendance}

\subsubsection*{Expectations}

\wss{What are your team's expectations regarding meeting attendance (being on
time, leaving early, missing meetings, etc.)?}

\subsubsection*{Acceptable Excuse}

\wss{What constitutes an acceptable excuse for missing a meeting or a deadline?
What types of excuses will not be considered acceptable?}

\subsubsection*{In Case of Emergency}

\wss{What process will team members follow if they have an emergency and cannot
attend a team meeting or complete their individual work promised for a team
deliverable?}

\subsection*{Accountability and Teamwork}

\subsubsection*{Quality} 

\wss{What are your team's expectations regarding the quality
of team members' preparation for team meetings and the quality of the
deliverables that members bring to the team?}

\subsubsection*{Attitude}

\wss{What are your team's expectations regarding team members' ideas,
interactions with the team, cooperation, attitudes, and anything else regarding
team member contributions?  Do you want to introduce a code of conduct?  Do you
want a conflict resolution plan?  Can adopt existing codes of conduct.}

\subsubsection*{Stay on Track}

\wss{What methods will be used to keep the team on track? How will your team
ensure that members contribute as expected to the team and that the team
performs as expected? How will your team reward members who do well and manage
members whose performance is below expectations?  What are the consequences for
someone not contributing their fair share?}

\wss{You may wish to use the project management metrics collected for the TA and
instructor for this.}

\wss{You can set target metrics for attendance, commits, etc.  What are the
consequences if someone doesn't hit their targets?  Do they need to bring the
coffee to the next team meeting?  Does the team need to make an appointment with
their TA, or the instructor?  Are there incentives for reaching targets early?}

\subsubsection*{Team Building}

\wss{How will you build team cohesion (fun time, group rituals, etc.)? }

\subsubsection*{Decision Making} 

\wss{How will you make decisions in your group? Consensus?  Vote? How will you
handle disagreements? }

\end{document}